\chapter{Especificaciones Técnicas de Medición}

Este anexo presenta las especificaciones del equipamiento principal utilizado para la caracterización experimental del acoplador de línea ramificada y los valores de referencia para evaluación del desempeño.

\section{Analizador Vectorial de Redes (VNA)}

\begin{table}[h]
\centering
\begin{tabular}{|l|l|}
\hline
\textbf{Parámetro} & \textbf{Especificación} \\
\hline
Modelo & Keysight E5071C ENA \\
Rango de frecuencias & 9 kHz a 8.5 GHz \\
Número de puertos & 4 puertos \\
Impedancia del sistema & 50 Ω (nominal) \\
Potencia de salida & -55 dBm a +10 dBm \\
Rango dinámico & >123 dB (10 MHz a 6 GHz) \\
Precisión de magnitud & $\pm$0.005 dB \\
Precisión de fase & $\pm$0.3° \\
\hline
\end{tabular}
\caption{Especificaciones del analizador vectorial de redes.}
\end{table}

\section{Kit de Calibración y Cables}

El kit de calibración utilizado es el Keysight 85052D 3.5mm (DC a 26.5 GHz), con una precisión ≤ 0.5\% hasta 8 GHz y un coeficiente de reflexión residual ≤ -40 dB. Los cables de medición son Gore PHASEFLEX con impedancia de 50 Ω ±1 Ω, pérdida de inserción ≤ 0.45 dB/m a 8 GHz y estabilidad de fase con flexión ≤ ±6° a 8 GHz.

\section{Valores de Referencia para Acopladores de Línea Ramificada}

\begin{table}[h]
\centering
\begin{tabular}{|l|c|c|c|c|}
\hline
\textbf{Parámetro} & \multicolumn{4}{c|}{\textbf{Frecuencia (GHz)}} \\
\cline{2-5}
& \textbf{4.0} & \textbf{5.0} & \textbf{6.0} & \textbf{8.0} \\
\hline
Acoplamiento (dB) & $3.0 \pm 0.3$ & $3.0 \pm 0.4$ & $3.0 \pm 0.5$ & $3.0 \pm 0.7$ \\
\hline
Directividad (dB) & $>25$ & $>23$ & $>20$ & $>18$ \\
\hline
Pérdidas de inserción (dB) & $<0.4$ & $<0.5$ & $<0.6$ & $<0.8$ \\
\hline
Desbalance de fase (°) & $90 \pm 2$ & $90 \pm 3$ & $90 \pm 4$ & $90 \pm 6$ \\
\hline
VSWR & $<1.2$ & $<1.3$ & $<1.4$ & $<1.6$ \\
\hline
\end{tabular}
\caption{Valores de referencia típicos para acopladores de línea ramificada comerciales de 3dB.}
\end{table}

\vspace{-0.5cm}
\section*{Certificación del Equipamiento}
\vspace{-0.3cm}
Todo el equipamiento utilizado fue verificado y calibrado según protocolos del fabricante. Los certificados de calibración están archivados en el laboratorio de Medidas Electrónicas de la UTN FRBA y cumplen con las normas ISO/IEC 17025:2017.